% =============================================================
%  ENLACE B: Canal 12 -- FCEFyN
%  Responsables: Arteaga Barrera, Baccino
%  Los números salen de datos/enlace_B.tex
% =============================================================
\renewcommand{\seccorta}{Enlace B}
\section{Enlace B: \EBnombre}

% ---------------------------------------------------------------
\begin{frame}{Definición del enlace y relevamiento geográfico}
  \small
  \begin{tabularx}{\textwidth}{@{}lLL@{}}
    \toprule
     & \textbf{Estación A: \EBunoNombre} & \textbf{Estación B: \EBdosNombre} \\ \midrule
    Ubicación & \EBunoUbic & \EBdosUbic \\
    Latitud / Longitud & \EBunoLat\ / \EBunoLon & \EBdosLat\ / \EBdosLon \\
    Cota del terreno (s.n.m.) & \EBunoCota & \EBdosCota \\
    Altura máxima disponible & \EBunoMax\ (torre) & \EBdosMax\ (mástil) \\
    Altura de instalación propuesta & \textbf{\EBunoAlt} & \textbf{\EBdosAlt} \\
    Cota total de la antena & \EBunoAntena & \EBdosAntena \\
    Azimut & \EBazUnoDos\ (A$\to$B) & \EBazDosUno\ (B$\to$A) \\
    Elevación (tilt) & \EBtiltUno\ (down-tilt) & \EBtiltDos\ (up-tilt) \\
    \bottomrule
  \end{tabularx}
  \vspace{6pt}
  \begin{columns}[T]
    \column{0.32\textwidth}
    \begin{block}{Distancia}\centering\Large\EBdist\end{block}
    \column{0.32\textwidth}
    \begin{block}{$\Delta$ terreno}\centering\Large\EBdifTerreno\end{block}
    \column{0.32\textwidth}
    \begin{block}{$\Delta$ antenas}\centering\Large\EBdifAntenas\end{block}
  \end{columns}
\end{frame}

% ---------------------------------------------------------------
\begin{frame}{Perfil del terreno y obstáculos}
  \begin{itemize}
    \item Perfil obtenido con \textbf{UISP Design Center} y validado en \textbf{Google Earth Pro}.
    \item Topografía tipo \textbf{``batea''}: ambas estaciones en zonas más altas que la depresión
          central (La Cañada y el centro de la ciudad).
    \item Los obstáculos críticos \textbf{no son naturales}: es el denso bloque de edificios de
          \textbf{Nueva Córdoba}, a $d_1=\EBdUno$ km de la Estación A.
    \item UISP \textbf{subestima los obstáculos urbanos} (sin datos LIDAR en Córdoba):
          al exportar el KMZ a Google Earth con edificios 3D aparecen torres que cortan la LOS
          topográfica $\Rightarrow$ se recalcularon las alturas por trigonometría.
  \end{itemize}
  \begin{alertblock}{Resultado}
    Con \EBunoAlt\ / \EBdosAlt\ la LOS pasa $\approx$\EBdespeje\ por encima de las terrazas más altas.
  \end{alertblock}
\end{frame}

% ---------------------------------------------------------------
\begin{frame}{Simulación en UISP Design Center}
  \centering
  \figura{B_uisp.jpg}{0.9\textwidth}\\[4pt]
  \small LiteBeam 5AC Gen2 en ambos extremos, canal de 40 MHz, ruido urbano de $-88$ dBm:
  señal esperada \EBprxSim, capacidad \EBcapCanal. UISP no tiene datos LIDAR en Córdoba.
\end{frame}

% ---------------------------------------------------------------
\begin{frame}{Línea de vista y primera zona de Fresnel}
  \begin{columns}[c]
    \column{0.55\textwidth}
    \small
    Obstáculo crítico: $d_1=\EBdUno$ km, $d_2=\EBdDos$ km, $D=$ \EBdist, $f=\EBfrecNum$ GHz
    \[
      R_1 = 17.32\sqrt{\frac{\EBdUno\times\EBdDos}{\EBfrecNum\times 6.69}} = \textbf{\EBrFresnel}
    \]
    \begin{itemize}
      \item Despeje mínimo (60\,\% de $R_1$): \textbf{\EBrSesenta}
      \item Despeje medido en Google Earth: \textbf{\EBdespeje}
      \item Despeje logrado: \textbf{\EBdespejePct} de $R_1$ \checkmark
      \item $\Rightarrow A_{obs} = 0$ dB (supera el 60\,\% recomendado por UIT)
    \end{itemize}
    \column{0.42\textwidth}
    \centering
    \figura{B_google_earth.jpg}{\linewidth}\\
    {\scriptsize Altura medida sobre el obstáculo (Google Earth)}
  \end{columns}
\end{frame}

% ---------------------------------------------------------------
\begin{frame}{Determinación de las alturas mínimas de antena}
  \begin{columns}[c]
    \column{0.55\textwidth}
    \small
    \begin{itemize}
      \item \textbf{\EBunoAlt\ (Est. A) / \EBdosAlt\ (Est. B)}: despeje de \EBdespejePct\ y
            señal esperada estable de \EBprxSim.
      \item \textbf{+5 m (30/25 m) y +10 m (35/30 m):} solo agregan margen ante futuras
            construcciones o grúas.
      \item \textbf{Conclusión:} ya se supera el 60\,\% recomendado, sin pérdidas por difracción; subir
            más no se justifica estructuralmente y complica el mantenimiento.
    \end{itemize}
    \column{0.42\textwidth}
    \centering
    \figura{B_senal_uisp.jpg}{\linewidth}\\
    {\scriptsize Señal esperada en UISP Design Center}
  \end{columns}
\end{frame}

% ---------------------------------------------------------------
\begin{frame}{Selección de frecuencia y marco regulatorio}
  \begin{block}{Banda elegida: \EBfrec}
    \begin{itemize}
      \item Res. 581-MM/2018 y Res. 4653/2019-ENACOM (STIC-UC, uso compartido sin licencia).
      \item En enlaces fijos PtP permite antenas de alta ganancia \textbf{sin reducir la potencia
            del transmisor}, a diferencia de las sub-bandas inferiores.
    \end{itemize}
  \end{block}
  \small
  \begin{tabularx}{\textwidth}{@{}lL@{}}
    \toprule
    Banda & 5725 -- 5850 MHz \\
    Condición de uso & Compartido, sin licencia (STIC-UC) \\
    Límites & 30 dBm conducidos / 36 dBm PIRE (sin reducción por ganancia en PtP) \\
    Operación del enlace & 7 dBm conducidos / 30 dBm PIRE \\
    Restricciones & No interferir a servicios primarios ni reclamar protección frente a ellos \\
    \bottomrule
  \end{tabularx}
\end{frame}

% ---------------------------------------------------------------
\begin{frame}{Selección del ancho de canal}
  \small
  UISP con ruido \textbf{Urban} ($-88$ dBm, referido a 20 MHz y escalado con $10\log B$), $P_{RX}=$ \EBprxSim:\\[4pt]
  \centering
  \begin{tabular}{@{}lcccc@{}}
    \toprule
    \textbf{Ancho de canal} & 10 MHz & 20 MHz & \cellcolor{naranja!25}\textbf{40 MHz} & 80 MHz \\ \midrule
    Capacidad (UISP) & 44.2 Mbps & 66.3 Mbps & \cellcolor{naranja!25}91.8 Mbps & 149 Mbps \\
    Modulación (UISP) & 6x 64QAM & 4x 16QAM & \cellcolor{naranja!25}4x 16QAM & 2x QPSK \\
    Ruido térmico $-174+10\log B$ & $-104$ & $-101$ & \cellcolor{naranja!25}$-98$ & $-95$ dBm \\
    Ruido urbano & $-91$ & $-88$ & \cellcolor{naranja!25}$-85$ & $-82$ dBm \\
    SNR & 21 dB & 18 dB & \cellcolor{naranja!25}15 dB & 12 dB \\
    Throughput ($\approx$70\,\%) & $\sim$31 & \textcolor{red}{$\sim$46} & \cellcolor{naranja!25}\textbf{$\sim$64} & $\sim$104 Mbps \\
    \bottomrule
  \end{tabular}
  \raggedright
  \vspace{6pt}
  \begin{alertblock}{Selección: \EBcanal}
    Único ancho que cumple los 50 Mbps con modulación estable: 20 MHz no alcanza
    ($\sim$46 Mbps) y 80 MHz cae a QPSK con SNR de 12 dB. El ruido urbano queda
    $\approx$13 dB sobre el térmico (figura de ruido + emisiones de la ciudad).
  \end{alertblock}
\end{frame}

% ---------------------------------------------------------------
\begin{frame}{Selección del equipamiento: \EBequipo}
  \small
  \begin{columns}[T]
    \column{0.52\textwidth}
    \begin{tabularx}{\linewidth}{@{}lL@{}}
      \toprule
      Modelo & LBE-5AC-Gen2 (ambos extremos) \\
      Banda & 5150 -- 5875 MHz \\
      Potencia máx. TX & 25 dBm \\
      Sensibilidad & $-96$ dBm (1x BPSK) a $-65$ dBm (8x 256QAM 5/6) \\
      Modulación & airMAX ac: BPSK\dots256QAM \\
      Anchos PtP & 10/20/30/40/50/60/80 MHz \\
      Throughput & $>$450 Mbps TCP/IP (80 MHz) \\
      Antena & Parabólica InnerFeed, \EBganancia, dual lineal \\
      Interfaz & 1 $\times$ Gigabit Ethernet \\
      Alimentación & PoE pasivo 24 V, 0.3 A; 7 W máx. \\
      Ambiente & $-40$ a 70\,°C; 5--95\,\% HR; viento 200 km/h \\
      \bottomrule
    \end{tabularx}
    \column{0.45\textwidth}
    \textbf{Justificación}
    \begin{itemize}
      \item \textbf{Distancia:} con 23 dBi, UISP estima \EBprxSim\ a \EBdist.
      \item \textbf{Capacidad:} soporta todos los anchos y 256QAM.
      \item \textbf{Frecuencia:} opera nativamente en 5 GHz.
      \item \textbf{Interferencias:} haz estrecho del reflector; UISP permite ajustar ruido urbano.
      \item \textbf{Margen:} 25 dBm TX y $-96$ dBm de sensibilidad.
      \item \textbf{Costo:} amplia disponibilidad comercial y precio accesible.
    \end{itemize}
  \end{columns}
\end{frame}

% ---------------------------------------------------------------
\begin{frame}{Capacidad requerida}
  \begin{itemize}
    \item Sin demanda estipulada: se fija \textbf{\EBcapReq} como mínimo.
    \item Suficiente para tráfico institucional: video comprimido (H.264/H.265),
          telefonía IP y acceso a bases de datos.
    \item Con \EBcanal\ y \textbf{ruido urbano}, UISP negocia 16QAM y una capacidad de
          \textbf{\EBcapCanal} $\Rightarrow$ throughput esperado \textbf{\EBcapRuido}.
  \end{itemize}
  \vspace{6pt}
  \begin{block}{Verificación en el peor escenario}
    \centering\large
    Throughput esperado ($\approx$64 Mbps) \;$>$\; Requerido (50 Mbps) \checkmark
  \end{block}
  \vspace{4pt}
  \small Se prioriza estabilidad y disponibilidad por sobre picos teóricos inalcanzables en
  espectro no licenciado.
\end{frame}

% ---------------------------------------------------------------
\begin{frame}{Atenuación de espacio libre y presupuesto de RF}
  \begin{columns}[T]
    \column{0.47\textwidth}
    \textbf{Atenuación de espacio libre}
    \begin{align*}
      A_0 &= 92.44 + 20\log(5.8) + 20\log(6.69)\\
          &= 92.44 + 15.27 + 16.51\\
          &= \textbf{124.22 dB}
    \end{align*}
    \small UISP no muestra $A_0$ explícitamente; solo la señal final esperada.
    \column{0.5\textwidth}
    \small
    \begin{tabular}{@{}lr@{}}
      \toprule
      $P_{TX}$ (criterio conservador) & \EBptx \\
      $L_{TX}$, $L_{RX}$ (InnerFeed) & 0 dB \\
      $G_{TX}$ & +23 dBi \\
      $A_0$ & $-124.2$ dB \\
      $A_{obs}$ (Fresnel despejada) & 0 dB \\
      $G_{RX}$ & +23 dBi \\ \midrule
      $P_{RX}$ \textbf{teórica} & \textbf{\EBprxTeo} \\
      $P_{RX}$ \textbf{simulada (UISP)} & \textbf{\EBprxSim} \\
      \bottomrule
    \end{tabular}
  \end{columns}
  \vspace{6pt}
  \begin{alertblock}{Discrepancia $\approx$1.2 dB}
    UISP probablemente calcula $A_0$ con otra frecuencia dentro de la banda
    (la teoría usa el extremo superior, 5.8 GHz).
  \end{alertblock}
\end{frame}

% ---------------------------------------------------------------
\begin{frame}{PIRE y verificación regulatoria}
  \[
    \text{PIRE} = P_{TX} + G_{TX} = 7\ \text{dBm} + 23\ \text{dBi} = \textbf{30 dBm (1 W)}
  \]
  \begin{columns}[T]
    \column{0.5\textwidth}
    \begin{block}{Res. 581-MM/2018, 5725--5850 MHz}
      \small
      \begin{itemize}
        \item Potencia conducida máx.: \textbf{30 dBm}
        \item PIRE máx.: \textbf{36 dBm}
        \item PtP fijo: antenas de alta ganancia sin reducir potencia
      \end{itemize}
      Con 7 dBm / 30 dBm el enlace cumple \textbf{holgadamente}.
    \end{block}
    \column{0.46\textwidth}
    \begin{alertblock}{¿Por qué 7 dBm?}
      \small
      UISP (modo Auto, Argentina) limita la PIRE a 30 dBm, el valor de las sub-bandas
      5250--5725 MHz. Se adopta como \textbf{criterio conservador}: en banda compartida
      conviene la mínima potencia que asegure el enlace, sin elevar el piso de ruido
      ni interferir a otros usuarios.
    \end{alertblock}
  \end{columns}
\end{frame}

% ---------------------------------------------------------------
\begin{frame}{Sensibilidad, modulación y margen de enlace}
  \begin{columns}[T]
    \column{0.5\textwidth}
    \small
    \begin{tabular}{@{}lcc@{}}
      \toprule
      \textbf{MCS} & $S_{RX}$ & $M$ ($P_{RX}=-70$) \\ \midrule
      8x 256QAM 3/4 & $-69$ dBm & \textcolor{red}{$-1$ dB} \\
      4x 16QAM 3/4  & $-86$ dBm & \textcolor{verde}{\textbf{16 dB}} \\
      1x BPSK 1/2   & $-96$ dBm & 26 dB \\
      \bottomrule
    \end{tabular}
    \[ M = P_{RX} - S_{RX} \]
    {\scriptsize Sensibilidades del fabricante con ruido térmico: en entorno urbano el margen
    efectivo es menor y manda el SNR.}
    \column{0.47\textwidth}
    \small
    \begin{itemize}
      \item \textbf{256QAM:} margen negativo, no sostenible.
      \item \textbf{16QAM:} margen de 16 dB por sensibilidad.
      \item \textbf{Rate Adaptation (airOS):} baja el MCS al achicarse el margen; la
            sensibilidad mejora hasta $-96$ dBm (BPSK) y el enlace \textbf{no se corta},
            a costa de menos bits/Hz y menor throughput.
    \end{itemize}
  \end{columns}
\end{frame}

% ---------------------------------------------------------------
\begin{frame}{Nivel de ruido y SNR}
  Ruido urbano de UISP escalado a 40 MHz: \EBruido
  \[
    \text{SNR} = P_{RX} - P_{ruido} = -70 - (-85) = \textbf{\EBsnr}
  \]
  \begin{itemize}
    \item Alcanza para 16QAM ($\approx$15--18 dB), coincide con el 4x de UISP;
          insuficiente para 256QAM ($>$32 dB).
    \item Está en el \textbf{límite inferior} de 16QAM: ante desvanecimientos podría bajar a
          QPSK $\Rightarrow$ se puede subir $P_{TX}$ dentro de la norma.
  \end{itemize}
  \vspace{4pt}
  \begin{center}
    \begin{tikzpicture}[every node/.style={draw=verde, fill=grisclaro,
        rounded corners, font=\small, minimum height=0.8cm, align=center}]
      \node (a) {Ruido $\uparrow$};
      \node[right=0.8cm of a] (b) {SNR $\downarrow$};
      \node[right=0.8cm of b] (c) {Símbolos no\\distinguibles};
      \node[right=0.8cm of c] (d) {MCS $\downarrow$};
      \node[right=0.8cm of d] (e) {Throughput $\downarrow$};
      \foreach \x/\y in {a/b,b/c,c/d,d/e} \draw[->,thick,naranja] (\x) -- (\y);
    \end{tikzpicture}
  \end{center}
\end{frame}

% ---------------------------------------------------------------
\begin{frame}{Pérdidas adicionales, interferencias y extremo crítico}
  \begin{columns}[T]
    \column{0.48\textwidth}
    \begin{block}{Pérdidas por obstáculos}
      \small
      Despeje real de \EBdespeje\ sobre \EBrFresnel\ (\EBdespejePct).
      Al superar el 60\,\% (\EBrSesenta), según las recomendaciones UIT
      $A_{obs}=0$ dB.
    \end{block}
    \begin{block}{Interferencias}
      \small
      En banda no licenciada no se puede garantizar ausencia de interferencias.
      En campo se usa \textbf{airView} (analizador de espectro de airOS) para elegir el
      canal con menor ruido.
    \end{block}
    \column{0.48\textwidth}
    \begin{alertblock}{Extremo más crítico: Est. B (FCEFyN, \EBdosAlt)}
      \small
      \begin{itemize}
        \item Cota más baja (\EBdosCota\ vs.\ \EBunoCota).
        \item Apunta al NNO: su lóbulo frontal abarca la densidad edilicia y las redes
              Wi-Fi de Nueva Córdoba y el Centro.
        \item Canal 12 está más elevado y en una zona residencial de menor densidad.
      \end{itemize}
    \end{alertblock}
  \end{columns}
\end{frame}

% ---------------------------------------------------------------
\begin{frame}{Orientación de las antenas}
  \begin{columns}[T]
    \column{0.42\textwidth}
    \begin{tabular}{@{}lcc@{}}
      \toprule
       & \textbf{Azimut} & \textbf{Tilt} \\ \midrule
      Est. A & \EBazUnoDos & \EBtiltUno \\
      Est. B & \EBazDosUno & \EBtiltDos \\
      \bottomrule
    \end{tabular}\\[6pt]
    Polarización: dual lineal (V y H).
    \column{0.55\textwidth}
    \small
    \begin{enumerate}
      \item \textbf{Alineación gruesa:} brújula o puntos notables de la ciudad.
      \item Activar la radio de gestión de \textbf{2.4 GHz} y conectar la app UNMS.
      \item \textbf{Alineación fina:} leer RSSI en vivo y ajustar la rótula de 3 ejes hasta el
            pico ($\approx$\EBprxSim) en \textbf{Chain 0 y Chain 1}.
      \item Ajustar los herrajes.
    \end{enumerate}
  \end{columns}
\end{frame}

% ---------------------------------------------------------------
\begin{frame}{Montaje, cableado, alimentación y protección}
  \begin{columns}[T]
    \column{0.48\textwidth}
    \begin{block}{Sistema de montaje}
      \small
      Kit de montaje a mástil incluido (ajuste independiente de azimut y elevación).
      \begin{itemize}
        \item \textbf{Est. A:} caño portante en la torre de celosía de ARTEAR.
        \item \textbf{Est. B:} mástil tubular arriostrado en el techo de FCEFyN (sobre LARFyM).
        \item 980 g con montaje; carga de viento 275 N a 200 km/h $\Rightarrow$ esfuerzo mínimo.
      \end{itemize}
    \end{block}
    \column{0.48\textwidth}
    \begin{block}{Cableado y protección}
      \small
      \begin{itemize}
        \item UTP/FTP Cat.\ 5e o 6, $\approx$\EBcable\ ($<$100 m de Ethernet).
        \item RJ45 blindados con contacto de masa.
        \item PoE pasivo 24 V, 0.3 A Gigabit (incluido).
        \item Protector \textbf{ETH-SP-G2} en la base del mástil, conectado a la PAT del edificio.
      \end{itemize}
    \end{block}
  \end{columns}
\end{frame}

% ---------------------------------------------------------------
\begin{frame}{Optimización del enlace}
  \begin{block}{Configuración propuesta}
    \EBequipo, \EBunoAlt/\EBdosAlt, canal de \EBcanal\ en 5.8 GHz, $P_{TX}=$ \EBptx
  \end{block}
  \small
  \begin{tabularx}{\textwidth}{@{}lL@{}}
    \toprule
    \textbf{Alternativa} & \textbf{Evaluación} \\ \midrule
    Canal de 80 MHz & Más capacidad teórica, pero con ruido urbano el SNR cae a 12 dB y queda en QPSK. \\
    Mayor altura & Más costo y exposición a descargas; el despeje de Fresnel ya está garantizado. \\
    Antenas de mayor ganancia & Encarece el proyecto y suma carga de viento; $-70$ dBm ya alcanza. \\
    \textbf{Mayor potencia} & \textbf{Mejora posible:} la norma permite subir de 7 a 25 dBm (máx.\ del equipo) si el ruido medido con airView supera lo estimado. \\
    \bottomrule
  \end{tabularx}
\end{frame}
